\section{树}

\subsection{重链剖分（LCA，两点距离，祖先判断、\texorpdfstring{$k$}{k} 级祖先）}

经过重链剖分，一棵树被拆成若干链，链内是重边，链外是轻边．
从根节点出发，每经过一次轻边，当前子树大小至少减半．
因此在任意路径上，轻边数量至多 $O(\log n)$ 条，
即路径总能拆成 $O(\log n)$ 段连续链区间．

关于 dfn 序（0-based）的性质：
\begin{itemize}[noitemsep]
    \item 一条链上自上而下，dfn 构成递增的连续区间，可据此在链上寻址、对路径操作；
    \item 子树所有节点的 dfn 构成的集合是连续区间，可据此对子树操作．
\end{itemize}

我们使用：
\begin{itemize}[noitemsep]
    \item {\tt top[u]} 表示 {\tt u} 所处链的头部；
    \item {\tt in[u]} 表示 {\tt u} 的 dfn；
    \item {\tt [in[u], out[u])} 表示 {\tt u} 子树的 dfn 区间；
    \item {\tt seq[i]} 将 dfn 序 {\tt i} 转化为对应节点编号．
\end{itemize}

节点编号是 0-based 的．

\lstinputlisting[caption={hld-base.cpp},label={hld-base},language=C++]{Chapter/tree/code/hld-base.cpp}

\subsection{重链剖分（换根后的父节点、子树大小、LCA）}

{\tt get\_rooted\_fa} 和 {\tt get\_rooted\_sz} 的复杂度都是 $O(\log \deg p)$．

\lstinputlisting[caption={hld-rooted.cpp},label={hld-rooted},language=C++]{Chapter/tree/code/hld-rooted.cpp}

\newpage

\subsection{重链剖分（线段树维护路径信息、子树信息）}

使用线段树维护 dfn 区间．

\lstinputlisting[caption={hld-sgt.cpp},label={hld-sgt},language=C++]{Chapter/tree/code/hld-sgt.cpp}

如果要维护有向的 dfn 区间（比如点权的结合（如仿射函数的复合）不满足交换律），
那么需要为 {\tt Info} 设计正反两向区间信息 {\tt Info::forw/backw}．
合并 {\tt Info} 时，{\bf 正向信息正向拼接，反向信息反向拼接}．

\lstinputlisting[caption={hld-sgt-directed.cpp},label={hld-sgt-directed},language=C++]{Chapter/tree/code/hld-sgt-directed.cpp}

\newpage

\subsection{直径}

经过子树根节点的直径由最深路径和次深路径组成，使用树形 DP 解决，支持负边权，时间复杂度为 $O(n)$．

\lstinputlisting[caption={diameter.cpp},label={diameter},language=C++]{Chapter/tree/code/diameter.cpp}

\subsection{重心}

去除重心后，每棵子树大小不超过 $n / 2$，据此寻找重心．
求有根树每棵子树的重心，时间复杂度为 $O(n)$．

\lstinputlisting[caption={centroid.cpp},label={centroid},language=C++]{Chapter/tree/code/centroid.cpp}

\newpage
